\documentclass[12pt,a4paper]{article}

\usepackage[T2A]{fontenc}
\usepackage[utf8]{inputenc}
\usepackage[english,russian]{babel}

\usepackage{geometry}
\geometry{margin=2.4cm}
\usepackage{microtype}
\usepackage{setspace}
\onehalfspacing
\emergencystretch=3em
\tolerance=2000

\usepackage{booktabs}
\usepackage{tabularx}
\usepackage{array}
\usepackage{ragged2e}
\usepackage{caption}
\newcolumntype{Y}{>{\RaggedRight\arraybackslash}X}
\newcolumntype{L}[1]{>{\RaggedRight\arraybackslash}p{#1}}
\captionsetup[table]{font=small, labelfont=bf}

\usepackage{enumitem}
\usepackage{amsmath}

\title{Пояснительная записка к решению тестового задания}
\author{Кузнецов Дмитрий}
\date{}

\begin{document}
\maketitle

\begin{abstract}
Задание просит две вещи: придумать, как встроить генеративный ИИ в квест-сценарий (Часть 2), и предложить метрики, которые скажут об участнике больше, чем прошёл / не прошёл (Часть 3). Моя основная идея простая: ИИ в квесте - это советник, который иногда ошибается, как и в жизни. Платформа оценивает не сами ответы, а то, как человек работает с советником: что спрашивает, проверяет ли рассуждения, когда доверяет, а когда нет - и что остаётся у него в голове, когда советник выключен.

В этом тексте я объясняю, почему выбрана именно такая схема, какие у неё риски и как каждая метрика считается из логов прохождения. Всё описанное реализовано в репозитории.
\end{abstract}

\pagebreak

\section{Разрешить нельзя запретить}

Начнём с главного вопроса: что вообще делать с ИИ на обучающей платформе? Вариантов два, и оба плохие. Запретить? Невозможно проконтролировать: участник откроет чат-бота на телефоне, и мы просто перестанем видеть, как он работает. Дать свободный доступ? Тут уже есть неприятные данные. В полевом эксперименте Бастани и коллег школьники с обычным GPT-4 решали тренировочные задачи на 48\% лучше - а потом, на контрольной без ИИ, показали результат на 17\% хуже тех, кто занимался вообще без модели \cite{bastani2025}. Модель сделала работу за них, знание не осталось. Показательно, что версия с ограждениями (подсказывает, но не выдаёт готовый ответ) этот вред убрала.

Есть и обратная сторона: правильно спроектированный ИИ-тьютор в исследовании Кестина дал вдвое больший прирост знаний, чем занятие с преподавателем в активном формате \cite{kestin2025}. Сделаем микровывод: дело не в том, есть ли у человека доступ к ИИ, а в том, как устроено взаимодействие.

Отсюда выводится политика из трёх слов, на которой построено всё решение. ИИ \textbf{ограничен}: советует, но не решает и не действует. ИИ \textbf{прозрачен}: каждое обращение к модели видно и записывается в лог. ИИ \textbf{оцениваем}: баллы получает не ответ, а процесс работы с советом. И одно важное дополнение: оценивание встроено в замкнутый цикл обучения - опыт \(\to\) обратная связь \(\to\) теория \(\to\) разбор \(\to\) проверка без ИИ \(\to\) повторное прохождение. Известно, что значительный вклад в обучение вносит не само упражнение, а его разбор \cite{fanning2007}. А проверка без ассистента - единственный честный способ узнать, чему человек научился, а что просто взял у модели \cite{bastani2025, koedinger2007}.

\section{Интеграция ИИ}

\subsection{Где и зачем}

ИИ-аналитик доступен на 6 из 13 развилок (я добавил ещё один этап): обеденный пик, невыход курьера, поломка печи, демпинг конкурента, вечерний пик с дождём и последний час смены. Логика выбора точек простая: это ситуации, в которых реальный операционный директор и правда открыл бы чат с моделью. 

Ключевой принцип: ИИ советует, но не решает. Он не переключает сцены, не выбирает действия, не трогает состояние квеста. Решение всегда принимает человек - и именно это мы измеряем.

Шаги на такой развилке открываются строго по очереди:

\begin{enumerate}[leftmargin=*]
    \item участник выбирает и фиксирует собственное предварительное решение - пока он этого не сделал, недоступны ни совет, ни итоговый выбор. После фиксации изменить намерение нельзя.
    \item формулирует запрос к модели своими словами и получает совет (одно обращение на развилку);
    \item по желанию раскрывает рассуждение модели и выносит вердикт: выдерживает проверку или есть противоречие;
    \item принимает итоговое решение - любое: совпадающее с намерением, с советом или ни с тем, ни с другим.
\end{enumerate}

 Участника честно предупреждают на старте: советы ИИ могут быть ошибочными, оценивается работа с ними.

\subsection{Что на входе и что на выходе}

\textbf{Вход} модели - три вещи: машинно-читаемое описание ситуации из графа сценария (что случилось, какие действия доступны), текущий вектор KPI (очередь, загрузка, заказы, прибыль, SLA, NPS, команда) и свободный текстовый запрос участника.

\textbf{Выход} - структурированный объект: рекомендуемое действие, численный прогноз влияния на KPI, развёрнутое рассуждение и разбор (последний показывается только после решения). 

Самая важная деталь интерфейса: кнопка «Проверить обоснование» показывает только рассуждение модели, без вердикта. Правильный ответ система не подсказывает - вердикт выносит человек, а узнаёт, был ли он прав, только после решения. Это когнитивная форсирующая функция: небольшое обязательное усилие, которое мешает пассивно доверять \cite{bucinca2021, bucinca2025}. В прототипе обращение к модели симулирует детерминированная заглушка - все участники получают одинаковые советы, иначе оценка нечестная. Живую модель (например, через API) можно подключить в ту же точку.

\subsection{Как совет влияет на ход квеста}

Напрямую - никак, и это принципиально. Ход квеста двигает не совет, а реакция человека на совет. Остальное влияние - через оценку (раздел 3) и через обучение: после каждого решения на ИИ-развилке участник получает разбор совета и своего вердикта, а до решения может открыть карточку теории для изучения. 

\subsection{Риски и что с ними делать}

\begin{enumerate}[leftmargin=*]
    \item \textbf{Галлюцинации.} Модель может уверенно выдумывать: порекомендовать действие, которого нет среди вариантов узла, нарисовать прогноз с числами «из воздуха», сослаться на факты, которых в ситуации нет. \textbf{Решение.} Технический уровень: RAG по банку информации о данном узле плюс промпт-инжиниринг. Посеянные ошибки можно устроить архитектурно: держать два банка контекстов - актуальный и полуактуальный - и в заранее выбранных узлах подавать модели второй, тогда неправильный совет получается контролируемо. 

    \item \textbf{Непредсказуемость.} Я бы не запускал живую LLM во время контролируемого тестирования данного подхода. В эксперименте использовал бы заглушки, а в заранее выбранных местах - обманки. 

    \item \textbf{Нерелевантный запрос.} Участник может спросить не по делу или попытаться увести модель от задачи. \textbf{Решение.} Фильтр релевантности по ключевым сущностям задачи (эвристика). Нерелевантный запрос получает 0 за постановку, но прохождение не блокирует. Либо поставить более умную ступень - фильтровать запрос по смысловой близости к задаче (cosine similarity по эмбеддингам либо сторонняя модель-судья, но такой подход может оказаться медленным)

\end{enumerate}

\section{Метрики}

\begin{table}[htbp]
\centering
\footnotesize
\setlength{\tabcolsep}{3pt}
\renewcommand{\arraystretch}{1.25}
\caption{Процессные метрики квеста}
\label{tab:metrics}
\begin{tabularx}{\textwidth}{@{}L{0.2\textwidth} Y Y@{}}
\toprule
\textbf{Метрика} &
\textbf{Что измеряет} &
\textbf{Как считается} \\
\midrule

Качество запроса (К1) &
Умеет ли человек ставить задачу ИИ &
Рубрика 0--4: цель или ограничение; ссылка на текущее состояние; просьба сравнить варианты; просьба обосновать. Нерелевантный запрос - 0 \\

Точность проверки (К2) &
Находит ли посеянные противоречия в рассуждениях ИИ &
Доля верных вердиктов «выдерживает / противоречие»; отдельно - доля ошибочных советов, отклонённых итоговым выбором \\

Калибровка опоры (К3) &
Когда человек доверяет совету, а когда нет &
RAIR и RSR: переключения к верному совету и устойчивость к ошибочному - с учётом зафиксированного намерения \\

Обучаемость (К4) &
Становятся ли решения лучше по ходу дня &
Наклон качества решений по шагам; доля исправлений сразу после неудачи \\ 

Перенос (К5а) &
Применяет ли человек принцип в новой ситуации &
Согласованность на похожих развилках внутри дня + отдельная развилка в конце без ИИ \\

Саморегуляция (К5б) &
Темп и осознанность &
Время на решение; обращения к теории; сигналы офлоадинга (быстро и без проверки - за ошибочным советом)  \\ 

\midrule
Композит (PQ) &
Общий процессный балл &
Взвешенная сумма компонент, 0--100, только по доступным данным \\
\bottomrule
\end{tabularx}
\end{table}

\subsection{Калибровка опоры: RAIR и RSR}

Наивный способ измерить доверие - посчитать, как часто человек следует советам. Он не работает: если участник и сам собирался делать то, что потом посоветовал ИИ, при чём тут доверие? Совпадение мнений - не опора. Поэтому мы фиксируем намерение до совета и считаем только осмысленные случаи - те, где намерение расходилось с советом \cite{schemmer2023}:
\[
\mathrm{RAIR}=\frac{\#\{\text{совет верный, намерение иное, итог}=\text{совет}\}}
                   {\#\{\text{совет верный, намерение иное}\}},
\]
\[
\mathrm{RSR}=\frac{\#\{\text{совет ошибочный, намерение иное, итог}\neq\text{совет}\}}
                  {\#\{\text{совет ошибочный, намерение иное}\}}.
\]
В двух словах: RAIR (Relative AI Reliance) - «меня переубедили, когда переубедить стоило», RSR (Relative Self-Reliance) - «я устоял, когда совет был плох». Хорошая калибровка - высокие оба. Итоговый индекс - среднее RAIR и RSR.

\subsection{Точность проверки}

Здесь три разных вопроса: открывал ли участник рассуждение (действие), верно ли судил о нём (качество) и отклонил ли ошибочный совет делом (поведение). Баллы даются только за вердикт. В выборке есть и верные советы, поэтому нельзя будет применить стратегию, что везде противоречие. 
\subsection{Обучаемость и перенос}

\textbf{Обучаемость (К4)} - наклон прямой, проведённой через качество решений (0/1/2) по порядку шагов, плюс доля случаев, когда после неудачного решения следующее было лучше. Грубая, но честная оценка траектории.

\textbf{Перенос (К5а)} - самая важная метрика, и вот почему. Квест дважды учит одному принципу: «вариативность - враг пропускной способности» (экспресс-меню на пике, сет-меню для корпоратива). Ближний перенос - применил ли участник принцип во второй раз. Но главная проверка - «утро следующего дня»: отдельная развилка в конце, где ИИ недоступен, а ситуация изоморфна вчерашним. Если человек весь день опирался на советника, а утром сам выбрал расширить меню перед пиком - день прошёл зря, что бы ни говорил P\&L. Ровно этот эффект - хорошо с ИИ, хуже без него - задокументирован в \cite{bastani2025}, поэтому проверка без ассистента встроена в сам сценарий.

\subsection{Композитный балл}

\[
\mathrm{PQ}=100\cdot\frac{\sum_k w_k m_k}{\sum_k w_k},
\]
где сумма берётся только по доступным компонентам. Веса (числа требуют калибровки, взяты произвольно): опора - 0,25; проверка - 0,20; запрос - 0,15; перенос - 0,20; динамика - 0,10; отлов ошибок - 0,10. Это принципиально: если участник ни разу не обратился к ИИ, у него нет данных об опоре и проверке - «нет данных» не равно нулю баллов. 

Сделаем микровывод: шесть метрик отвечают на шесть разных вопросов - как человек спрашивает, как проверяет, когда доверяет, учится ли, переносит ли и как думает. Бинарная оценка не отвечает ни на один.

\subsection{Ограничения и план проверки}

\textbf{Дизайн стоит проверить экспериментом.} Два простых A/B: (а) готовый вердикт против собственного вердикта - действительно ли обязательное усилие улучшает калибровку \cite{bucinca2025}; (б) наличие разбора дня против его отсутствия - влияет ли дебрифинг на выбор в «утре следующего дня».

\section{Что мы получили?}

Платформа, которая тренирует и измеряет одно и то же: постановку задачи, проверку чужих рассуждений, калиброванное доверие, обучаемость и перенос. ИИ в ней - не бог и не запретный плод, а рабочий инструмент с известным свойством: иногда он ошибается. Главный вопрос, который задаёт участнику такой квест, - не «угадал ли ты правильную кнопку», а «что останется, когда ИИ завтра выключат?» Именно это квест и проверяет последней развилкой. 

\begin{thebibliography}{9}

\bibitem{bastani2025} Bastani H., Bastani O., Sungu A., Ge H., Kabakc{\i} {\"O}., Mariman R. \textit{Generative AI without guardrails can harm learning: Evidence from high school mathematics}. PNAS 122(26), e2422633122, 2025.

\bibitem{kestin2025} Kestin G., Miller K., Klales A., Milbourne T., Ponti G. \textit{AI tutoring outperforms in-class active learning: an RCT introducing a novel research-based design in an authentic educational setting}. Scientific Reports 15, 17458, 2025.

\bibitem{bucinca2021} Bu{\c{c}}inca Z., Malaya M.B., Gajos K.Z. \textit{To Trust or to Think: Cognitive Forcing Functions Can Reduce Overreliance on AI in AI-assisted Decision-making}. PACM HCI 5(CSCW1), 188, 2021.

\bibitem{bucinca2025} \textit{Cognitive Forcing for Better Decision-Making: Reducing Overreliance on AI Systems Through Partial Explanations}. PACM HCI (CSCW), 2025. doi:10.1145/3710946.

\bibitem{schemmer2023} Schemmer M., K{\"u}hl N., Benz C., Bartos A., Satzger G. \textit{Appropriate Reliance on AI Advice: Conceptualization and the Effect of Explanations}. ACM IUI, 2023.

\bibitem{gerlich2025} Gerlich M. \textit{AI Tools in Society: Impacts on Cognitive Offloading and the Future of Critical Thinking}. Societies 15(1), 6, 2025.

\bibitem{fan2025} Fan Y. et al. \textit{Beware of metacognitive laziness: Effects of generative AI on learning motivation, processes, and performance}. British Journal of Educational Technology 56(2), 2025.

\bibitem{koedinger2007} Koedinger K.R., Aleven V. \textit{Exploring the Assistance Dilemma in Experiments with Cognitive Tutors}. Educational Psychology Review 19, 2007.

\bibitem{chi2014} Chi M.T.H., Wylie R. \textit{The ICAP Framework: Linking Cognitive Engagement to Active Learning Outcomes}. Educational Psychologist 49(4), 2014.

\bibitem{fanning2007} Fanning R.M., Gaba D.M. \textit{The Role of Debriefing in Simulation-Based Learning}. Simulation in Healthcare 2(2), 2007.

\bibitem{unesco2024} UNESCO. \textit{AI competency framework for students}. UNESCO, 2024.

\bibitem{ailit2026} European Commission, OECD. \textit{Empowering Learners for the Age of AI: AI Literacy Framework for Primary and Secondary Education}. Final version, 2026.

\end{thebibliography}

\end{document}
